% ===================================================================
%  TABELO N-RO 7 — La vilaĝo dum la vintro. — La kampara domo dum la
%  printempo.
%  Traduction 2026 d'apres texto/io, controlee sur texto/fr et
%  texto/es. Consigne : voir texto/eo/10-tabelo-01.tex.
% ===================================================================

%%K t07-noto-1 noto
\VUnotes{143.2mm}{%
(*) Vidu pri la manĝaj detaloj, la tabelojn 6 kaj 13.%
}

%%K t07-apar-1 apar
\VUsaut{4.0mm}
\VUcentre{13.2pt}{90}{DUA SERIO}
\VUsaut{3.0mm}
\VUfilet{16mm}
\VUsaut{5.6mm}
\VUtitre{15.0pt}{60}{TABELO N\textsuperscript{o} 7}
\VUsaut{3.0mm}

%%K t07-tit-1 sub
\VUcentre{12.0pt}{0}{\textit{Unua sceno.}}
\VUsaut{3.0mm}

%%K t07-tit-2 sub
\VUpk{10.2pt}{40}{La vilaĝo dum la vintro.}
\VUsaut{4.0mm}

%%K t07-c1-01-1 p
1. --- Dum la novjaraj ferioj, mi akompanis mian patron en Novurbon, vilaĝeton
kie li bezonis trakti kelkajn komercajn aferojn.

%%K t07-c1-02-1 p
2. --- Ni uzis la \VUgras{diliĝencon} \textsuperscript{(11)} kiu funkcias inter
la urbo kaj tiu burgo; sed ĉar estis tre malvarme, tiu vojaĝo en veturilo
malmulte komforta ne estis tre agrabla.

%%K t07-c1-03-1 p
3. --- Pro tio ni sentis veran plezuron kiam ni vidis la \VUgras{kondukiston}
haltigi sian veturilon sur la placo, antaŭ la \VUgras{gastejo}
\textsuperscript{(2)} de la \textit{Blanka Urso}. \VUgras{La gastejestro}
\textsuperscript{(4)} estis atendanta nin sur la sojlo de sia domo, kviete
fumante sian \VUgras{pipon}.

%%K t07-c1-03-2 p
Li invitis nin eniri kaj mi ne lasis min peti por instali min antaŭ la
sarmenta fajro kiu krakis en la fajrujo. Tiam mi tre mokis la akran
\VUgras{nordan venton} kiu grincigas la malnovan \VUgras{ŝildon}
\textsuperscript{(3)} kaj forportis per nigraj kirloj \VUgras{la fumon}
\textsuperscript{(1)} elirantan el la kameno.

%%K t07-c1-04-1 p
4. --- Estis la horo de la tagmanĝo. Sen pli longe atendi, ni bone manĝis la
simplan sed bongustan manĝon kiun oni servis al ni (*).

%%K t07-tit-3 sub
\VUpk{10.2pt}{40}{Kelkaj vizitoj.}
\VUsaut{3.0mm}

%%K t07-c1-05-1 p
5. --- Post la tagmanĝo, mi akompanis mian patron, kiu estas asekuristo, ĉe lia
klientaro. Unue ni vizitis la \VUgras{forĝiston} \textsuperscript{(5)}, kiu
loĝas apud la gastejo. Li estis okupata hufumi ĉevalon. Mi admiris la
fortikecon de lia kunulo kiu solide tenis sur sia leda antaŭtuko la
\VUgras{hufon} de la ĉevalo, dum la forĝisto fiksis la \VUgras{feron}
\textsuperscript{(6)} per fortaj martelbatoj.

%%K t07-c1-06-1 p
6. --- Poste ni iris al la \VUgras{ŝuisto} \textsuperscript{(7)} de la vilaĝo,
la maljuna Jakobo. Kvankam li havas kiel ŝildon belegan \VUgras{boton}, li
kontentiĝis per resuoli malnovan paron da ŝuoj truigita.

%%K t07-c1-06-2 p
La maljunulo diris al ni ridante : « En la urbo, mi estis “ŝufaristo” kaj mi
ne havis klientojn. Ĉi tie, mi estas “ŝuriparisto” kaj la laboro estas multa. »

%%K t07-c1-07-1 p
7. --- Li parolis al ni pri sia najbaro la \VUgras{razisto}
\textsuperscript{(8)}, kies klientaro ne estas kontenta. Liaj \VUgras{raziloj}
ĉiam estas breĉigitaj kaj liaj \VUgras{tondiloj} neniam tondas bone. Sed ĉar li
estas la sola razisto de la vilaĝo, oni estas devigita kompreneble razigi sin
kaj tondigi sian hararon de li. Cetere li estas avara kaj malafabla viro.

%%K t07-c1-08-1 p
8. --- Feliĉe, la ceteraj \VUgras{najbaroj} ne similas lin. Ekzemple la
\VUgras{ĉarfaristo} \textsuperscript{(9)} estas bonkora viro, laborema kaj
komplezema. Estas plezuro riparigi veturilon de li. Lia laboraĵo estas ĉiam
finita.

%%K t07-c1-09-1 p
9. --- Maljuna Jakobo ankaŭ ne plendis pri la \VUgras{selfaristo}
\textsuperscript{(10)}, kiun li kelkfoje helpas kudri la \VUgras{jungilarojn}
kiam urĝas la laboro.

%%K t07-c1-09-2 p
Mia patro demandis lin pri lia familio : « Ĉiuj fartas bone. Mia nepino estas
en la \VUgras{lernejo} \textsuperscript{(12)}. Ŝia \VUgras{instruistino}
\textsuperscript{(13)} estas tre kontenta pri ŝi, ŝi estas bona
\VUgras{lernanto} \textsuperscript{(14)}. Ŝi ne similas mian malbonan
knabaĉon filon; li pensas nur pri ludi. La \VUgras{instruisto}
\textsuperscript{(15)} ne sukcesas lernigi ion ajn de li. »

%%K t07-tit-4 sub
\VUsaut{3.0mm}
\VUpk{10.2pt}{40}{Vilaĝaj babiladoj.}
\VUsaut{3.0mm}

%%K t07-c1-10-1 p
10. --- La maljunulo poste rakontis al ni la novaĵojn de la vilaĝo. Oni estas
konstruonta novan lernejon, ĉar la instalita en la \VUgras{komunuma domo}
fariĝis tro malgranda. Cetere tio estas facila, ĉar sufiĉos aĉeti parton de la
ĝardeno de la \VUgras{muelisto}. Tio estos tre profita por la kompatinda viro.
Pro la malvarmo, la \VUgras{muelŝtonoj} de la \VUgras{muelejo}
\textsuperscript{(16)} ne plu povas mueli la tritikon, ĉar la \VUgras{rado}
\textsuperscript{(17)} estis frostige haltigita de la \VUgras{glacio}
\textsuperscript{(25)} de la \VUgras{rivereto} \textsuperscript{(23)}.

%%K t07-c1-11-1 p
11. --- « Pri konstruaĵo, ĉu vi vidis nian novan \VUgras{preĝejon}
\textsuperscript{(18)}? Ĝi estas tre pentrinda sub sia neĝa mantelo. La
\VUgras{korvoj} \textsuperscript{(19)}, kiuj ne plu rekonas sian malnovan
sonorilejon, triste sidas sur la granda arbo de la \VUgras{tombejo}
\textsuperscript{(20)} ».

%%K t07-c1-12-1 p
12. --- Tiu ideo pri la tombejo rememorigis al la ŝufaristo la personojn kiujn
mia patro estis koninta kaj kiuj mortis lastjare. « Kiom da \VUgras{tomboj}
\textsuperscript{(21)}, mia bona sinjoro, aldoniĝis al la ceteraj, sub la
\VUgras{cipreso} \textsuperscript{(22)}! La morto falĉis nian malnovan notarion.
La tutaj vilaĝanoj ĉeestis lian \VUgras{funebron}. »

%%K t07-c1-13-1 p
13. --- « Jen lia posteulo kiu glitkuras sur la rivereto. Li ĵus venis riparigi
de mi la rimenon de sia \VUgras{glitilo} \textsuperscript{(26)} kiu estis
rompita. »

%%K t07-c1-14-1 p
14. --- « Jen ankaŭ sur la bordo de la rivereto, la nova \VUgras{kampara
gardisto} \textsuperscript{(27)}. Li estas eksĝendarmo kiu terurigas la
bubaĉojn de la vilaĝo. Ili forkuras tuj kiam ili ekvidas liajn \VUgras{gamaŝojn}
\textsuperscript{(29)} kaj lian \VUgras{kepon} \textsuperscript{(28)}. »

%%K t07-c1-15-1 p
15. --- « Ha! jen la maljuna Jozefo kiu kondukas \VUgras{ĉareton da lignaj
faskoj} \textsuperscript{(30)} por la panisto. --- Tio estas vera, diris mia
patro, mi rekonas lian jungaĵon el \VUgras{muloj} \textsuperscript{(31)}. Mi ja
bezonas paroli al li, mi profitos la okazon. Ĝis revido, patro Jakobo. »

%%K t07-c1-16-1 p
16. --- Ĝuste kiam ni eliris, la knaboj de la vilaĝo forlasis la lernejon kaj
kurante ĵetis sin al la \VUgras{glitejo} \textsuperscript{(33)} kun risko fali
kaj rompi al si membron en la \VUgras{falo} \textsuperscript{(34)}. Unu el ili
prenis sian \VUgras{glitveturilon} \textsuperscript{(32)} ĉe la ĉarfaristo,
sidigis en ĝin unu el siaj kamaradoj, kaj foriris kurante.

%%K t07-c1-17-1 p
17. --- Aliaj ĵetis sin al \VUgras{neĝamaso} \textsuperscript{(37)} apud la
\VUgras{ponto} \textsuperscript{(40)} kaj ekbombardis la \VUgras{neĝan statuon}
\textsuperscript{(39)} kiun ili starigis dum la hieraŭa tago kaj al kiu ili
metis ĉapelon, pipon kaj ŝultran lernosakon.

%%K t07-c1-18-1 p
18. --- Ili malmulte zorgas pri la frostiga vento kaj la malvarmo. La plimulto
eĉ ne havas \VUgras{koltukegon} \textsuperscript{(35)} kaj, por revarmigi sin
post la batalo per \VUgras{neĝbuloj} \textsuperscript{(38)}, sufiĉas al ili
manpleno da tre varmaj \VUgras{kaŝtanoj} \textsuperscript{(42)} aĉetitaj de la
maljuna \VUgras{vendistino} Jakobina, kiu tremas pro malvarmo sub sia
\VUgras{ŝalaĉo} \textsuperscript{(43)} tro maldika.

%%K t07-tit-5 sub
\VUsaut{3.0mm}
\VUcentre{12.0pt}{0}{\textit{Dua sceno.}}
\VUsaut{3.0mm}

%%K t07-tit-6 sub
\VUpk{10.2pt}{40}{La kampara domo dum la printempo.}
\VUsaut{4.0mm}

%%K t07-c2-01-1 p
1. --- Malgraŭ ĉiuj plezuroj de la vintro, ni kun profunda ĝojo salutas la
revenon de la \VUgras{printempo}.

%%K t07-c2-01-2 p
Kiel gajigan tabelon aspektas farma korto, vespere, en la majo-monato!

%%K t07-c2-02-1 p
2. --- Ĉe la horizonto la \VUgras{suno} \textsuperscript{(1)} malrapide kuŝigas
sin, inundante la \VUgras{kamparon} \textsuperscript{(3)} per siaj purpuraj
\VUgras{radioj} \textsuperscript{(2)}. La diligenta \VUgras{plugisto}
\textsuperscript{(6)} rapidas plugi sian \VUgras{kampon} \textsuperscript{(4)}.

%%K t07-c2-02-2 p
Per sia \VUgras{plugilo} \textsuperscript{(5)} jungita al vigora
\VUgras{ĉevalo} \textsuperscript{(7)}, li tiras la \VUgras{sulkon}
\textsuperscript{(8)} en kiun la \VUgras{semanto} \textsuperscript{(9)} per
plena manpleno ĵetos la \VUgras{semon} kiu produktos la orajn spikojn.

%%K t07-c2-03-1 p
3. --- La \VUgras{falĉistoj} \textsuperscript{(11)} provizitaj per siaj
falĉiloj falĉis la herbon de la herbejo, la fojnfaristoj sekigis la
\VUgras{fojnon} \textsuperscript{(10)} turnante ĝin per \VUgras{forkoj}
\textsuperscript{(12)} kaj rastiloj, kaj jen ke ili revenas.

%%K t07-c2-03-2 p
Unuj marŝas antaŭ la peza veturilo tirata de bovoj; ili kunportas sian ilon
sur la ŝultro. Aliaj, pli maldiligentaj, komforte sin instalis sur la odoranta
fojno.

%%K t07-c2-04-1 p
4. --- Pasante ili faras amikan saluton al la \VUgras{ĝardenisto}
\textsuperscript{(16)} okupata en la \VUgras{fruktoĝardeno}
\textsuperscript{(13)} tranĉi la \VUgras{fruktarbojn} kaj grefti la
\VUgras{pirarbojn} \textsuperscript{(14)}. La \VUgras{prunarboj}
\textsuperscript{(15)} ekfloras, kaj, en la \VUgras{legomĝardeno}
\textsuperscript{(17)} bone sterkita, la legomoj, brasikoj kaj salatoj jam
bonfartas.

%%K t07-c2-04-2 p
Sur la mureto, bela \VUgras{pavo} \textsuperscript{(18)} etendas sian voston,
por admirigi ĝiajn belegajn plumojn.

%%K t07-tit-7 sub
\VUpk{10.2pt}{40}{Farmokorto.}
\VUsaut{3.0mm}

%%K t07-c2-05-1 p
5. --- La pordoj de la \VUgras{ĉevalejo} \textsuperscript{(19)} estas
malfermitaj kaj oni vidas la mangujon kaj la \VUgras{sternaĵon}
\textsuperscript{(23)} de la \VUgras{ĉevalino} \textsuperscript{(21)} kiun la
\VUgras{ĉevalservisto} \textsuperscript{(20)} ĵus maljungis. La
\VUgras{ĉevalido} \textsuperscript{(22)} tiel komika kun siaj longaj kruroj,
sekvas sian patrinon saltetante.

%%K t07-c2-06-1 p
6. --- Apud la \VUgras{puto} \textsuperscript{(29)} la \VUgras{bovinoj}
\textsuperscript{(25)} iras trinki en la ligna \VUgras{trogo}
\textsuperscript{(26)} kiu uzatas por ili kiel trinkejo. La \VUgras{bovido}
\textsuperscript{(24)} profitas la okazon por mamsuĉi, kaj la \VUgras{virbovo}
\textsuperscript{(27)} flarante la bonan odoron de la furaĝejo, ĝoje muĝas kaj
fiere levas la kapon kun \VUgras{kornoj} \textsuperscript{(28)} potencaj.

%%K t07-c2-07-1 p
7. --- Ĉiuj laboras. La \VUgras{servistino} \textsuperscript{(32)} mane tenas
la \VUgras{ŝnuron} \textsuperscript{(31)} de la puto. Ŝi ĉerpas akvon per
\VUgras{sitelo} \textsuperscript{(30)} por trinkigi la brutaron. Apud la
\VUgras{garbejo} \textsuperscript{(33)}, viro rastas la pajlaĵojn kaj antaŭ la
pordo de la \VUgras{hejmo} \textsuperscript{(34)} maljuna \VUgras{ŝpinistino}
\textsuperscript{(35)} per sia ŝpinrado kaj sia \VUgras{ŝpinstango}, ŝpinas la
\VUgras{linon} aŭ la \VUgras{kanabon} kiel en la tempo antaŭa.

%%K t07-tit-8 sub
\VUsaut{3.0mm}
\VUpk{10.2pt}{40}{La farmisto.}
\VUsaut{3.0mm}

%%K t07-c2-08-1 p
8. --- La \VUgras{farmisto} \textsuperscript{(36)} direktas tiujn ĉiujn
laborojn kaj donas instrukciojn al siaj servistoj. Li rekomendas subtegmentigi
la \VUgras{erpilon} \textsuperscript{(40)} kaj la \VUgras{pioĉon}
\textsuperscript{(39)} kiujn oni forgesis apud la \VUgras{kabano}
\textsuperscript{(38)} de la \VUgras{hundo} \textsuperscript{(37)}.

%%K t07-c2-09-1 p
9. --- Liaj krioj timigas \VUgras{kolombon} \textsuperscript{(41)}, kiu per
ĉiuj flugiloj flugas en la \VUgras{kolombejon} \textsuperscript{(42)}, kaj la
\VUgras{kokinojn} \textsuperscript{(43)} kiuj rapidas reeniri la
\VUgras{kokejon} \textsuperscript{(44)}.

%%K t07-c2-10-1 p
10. --- Antaŭ la \VUgras{ŝafejo} \textsuperscript{(48)}, jen la maljuna
\VUgras{paŝtisto} \textsuperscript{(45)} kiu rekondukas sian (ŝaf)\VUgras{aron}
\textsuperscript{(49)}. Tenante sian (paŝt)\VUgras{bastonon}
\textsuperscript{(46)}, kiu neniam mankas al li, same kiel lia \VUgras{bluzo}
\textsuperscript{(47)} senkolorigita, li kondukas al la stalo
\VUgras{virŝafojn} \textsuperscript{(50)}, \VUgras{ŝafinojn}
\textsuperscript{(53)}, \VUgras{ŝafidojn} \textsuperscript{(54)},
\VUgras{kaprinojn} \textsuperscript{(51)} kaj \VUgras{kapridojn}
\textsuperscript{(52)}. Lia fidela \VUgras{hundo} \textsuperscript{(55)} Arguso
laste venas kaj instigas per siaj bojoj la diligentecon de la malfruiĝantoj.

%%K t07-c2-11-1 p
11. --- Tiu tuta bruado kaj movado ne malhelpas la kvieton de la bona
\VUgras{laktodona bovino} \textsuperscript{(56)} kiu tintigas sian
\VUgras{sonorileton} \textsuperscript{(57)}, dum oni melkas ŝin antaŭ la pordo
de la \VUgras{stalo} \textsuperscript{(58)}.

%%K t07-c2-12-1 p
12. Ŝi ŝajnas kompreni, ke ŝi devas doni sian lakton, por ke la
\VUgras{farmistino} \textsuperscript{(60)} povu pretigi \VUgras{buteron} en la
\VUgras{buterfarilo} \textsuperscript{(61)} aŭ la bonegajn \VUgras{fromaĝojn}
\textsuperscript{(64)}, kiuj gutas de la breto metita sur la \VUgras{kuvego}
\textsuperscript{(63)}.

%%K t07-c2-13-1 p
13. --- La tuta \VUgras{laktejo} \textsuperscript{(59)} estas tenata tre pura de
la \VUgras{farmoservisto} \textsuperscript{(65)} kiu ĵus lavis la
\VUgras{laktopotojn} \textsuperscript{(62)} en la granda sitelo kiun li portas
per la mano.

%%K t07-tit-9 sub
\VUsaut{3.0mm}
\VUpk{10.2pt}{40}{La kortbirdejo.}
\VUsaut{3.0mm}

%%K t07-c2-14-1 p
14. --- Kun siaj dikaj lignoŝuoj, li marŝas iom mallerte, kaj tiel forkurigas
la \VUgras{meleagron} \textsuperscript{(68)}, la \VUgras{anasinojn}
\textsuperscript{(66)}, la \VUgras{viranasojn} \textsuperscript{(67)} kaj la
\VUgras{anserojn} \textsuperscript{(69)} kiuj large malfermas \VUgras{sian
bekon} \textsuperscript{(70)} kaj malkviete kriegas.

%%K t07-c2-14-2 p
La \VUgras{virkoko} \textsuperscript{(71)}, pli batalema, starigas sian ruĝan
\VUgras{kreston} \textsuperscript{(72)}, skuas la plumojn de sia \VUgras{vosto}
\textsuperscript{(73)} kaj aŭdigas « kokeriko » sonoregan.

%%K t07-c2-15-1 p
15. --- \VUgras{Kokin-patrino} \textsuperscript{(74)} rabas sur la
\VUgras{sterko} \textsuperscript{(81)} kun siaj \VUgras{(kok)idetoj}
\textsuperscript{(75)}. La \VUgras{porkido} \textsuperscript{(78)} forkuras al
sia patrino, la enorma \VUgras{porkino} \textsuperscript{(77)} kiun ni vidas
apud la porkejo.

%%K t07-c2-16-1 p
16. --- En siaj fakoj, la \VUgras{kunikloj} \textsuperscript{(79)} manĝas avide
la \VUgras{herbon} \textsuperscript{(80)} delikatan kiun alportas al ili la
filino de la farmisto. Johanino tre amas siajn kuniklojn kaj, ĉiutage,
serĉinte la freŝajn \VUgras{ovojn} \textsuperscript{(76)} en la kokejo, ŝi
venas viziti siajn amiketojn.
\VUsaut{6.0mm}
\VUfilet{22mm}
